\section{Introduction}\label{s:intro}

\subsection{The question and the kind of answer}

The question to which EXI-01-a is addressed is the oldest question of its discipline and the one most often thought to have no answer: what is the meaning of life? It has been asked in three ways, and an answer to it must say which way it takes it.

On the first, the question asks for the point of human existence as such: what human beings, or the universe that contains them, are for. Answers of this kind assign a purpose to the species or to the whole, and they are typically theological. The individual life is meaningful, on such answers, to the extent that it plays its part in that purpose. On the second, the question asks what makes an individual life meaningful: what a life must contain, or how it must be lived, for its subject to have lived a meaningful life rather than a wasted, empty or trivial one. \citet{Metz2013} calls this the question of meaning \emph{in} life, as against the meaning \emph{of} life, and most analytic work since the 1970s addresses it. On the third, the question is conceptual: what we are asking when we ask whether a life has meaning, and what distinguishes that question from the questions whether the life was happy or whether it was morally good. \citet{Metz2001} argued that the concept is a cluster, united by a family of ideas: purposes worth pursuing, transcending one's animal nature, being an apt object of esteem or admiration, and making a narrative whole of one's life.

EXI-01-a takes the question in the second way, and its title keeps the wording of the first. That choice is examined at \Lm{1} and again at \Lm{9.1}. The answer it gives belongs to a family of answers that \citet{Metz2013} classifies as naturalist, since they do not make meaning depend on a spiritual realm, and within that family to the subclass of hybrid views, which make meaning depend both on the subject's attitudes and on something that does not depend on them. Its source, acknowledged in the text, is the fitting-fulfilment view of \citet{Wolf1997,Wolf2010}, on which meaning arises when the subject's active engagement meets an object worthy of it. Its distinctive additions are three: it makes the degree of meaning explicit, it isolates the assumption that there are objects of the required kind, and it claims that meaning so understood is available to every agent.

\subsection{The debate before EXI-01-a}

A commentary must place its text in the debate from which it comes. We give the history in outline here; the lemmata at which each strand becomes relevant return to it in detail.

\paragraph{Aristotle.} The \emph{Nicomachean Ethics} does not ask what the meaning of life is, but it asks the closest ancient question: what the chief good of a human life is, the end for whose sake everything else is done. Its answer, that the human good is activity of the soul in accordance with excellence over a complete life (\emph{NE} I.7, 1097b22--1098a20), contains three ideas that recur in EXI-01-a: that the good of a life lies in an activity (\emph{energeia}) rather than a state or a possession, that the activity must be of something worth doing, since the excellence in question is the excellence of the function proper to human beings, and that the assessment is of a life as a whole and not of a moment (\emph{NE} I.7, 1098a18--20; I.10, 1100a10--1101a21). The first idea is the ancestor of \enquote{active engagement} in \Lm{2.5}, the second of \enquote{objective worth} in \Lm{2.7}, and the third of \enquote{a life} in \Lm{2.1}. Aristotle also shows how far apart the ancient and modern questions are: his question concerns happiness (\emph{eudaimonia}), and EXI-07 argues that happiness is not the goal of life and is best understood as a by-product of engagement in worthwhile activities. EXI-01-a separates meaning from happiness in a way Aristotle does not.

\paragraph{Tolstoy and Camus.} The modern question is the question of the crisis. Tolstoy, in the \emph{Confession} \citep{Tolstoy1987}, describes how, at the height of his success, his life came to a halt before questions he could not answer: what would come of what he was doing, and whether there was any meaning in his life that would not be destroyed by the death that awaited him. He compared his condition to that of a traveller in an Eastern fable who, fleeing a beast, hangs from a branch over a well in which a dragon waits, while two mice gnaw at the branch, and who licks drops of honey from the leaves. The honey, the goods of his life, no longer tasted sweet, because he could see the dragon and the mice. \citet{Camus1955} opened \emph{The Myth of Sisyphus} with the claim that whether life is worth living is the one serious philosophical problem, and located the absurd in the confrontation between the human demand for meaning and the silence of the world. Both texts present the question of meaning as a question that arises when the subject steps back from her life; both take the answer to depend on whether something outside the life could vindicate it. EXI-01-a's Remark~5 engages that structure directly, and Section~\ref{s:absurd} examines it.

\paragraph{Taylor, Nagel and Wiggins.} The analytic discussion begins with three texts of the 1970s. \citet[ch.~18]{Taylor1970} took the myth of Sisyphus as his model of a meaningless existence and asked what would make it meaningful. Giving Sisyphus a purpose outside his labor, a temple to be built from the stones, would not do, since the completed temple would only be a scene of boredom; implanting in him a desire to roll stones would, since he would then be doing what he most wanted to do. Taylor concluded that meaning comes from within, from the will to do what one does. \citet{Nagel1971} argued that the sense of the absurd arises from a collision within the reflective subject between the seriousness with which we take our lives and the perpetual possibility of regarding everything about which we are serious as arbitrary, and that the appropriate response is irony. \citet{Wiggins1976} argued against Taylor that a purely subjective account of meaning cannot be lived from the inside, since the agent who finds her life meaningful takes its objects to be worthy of her concern, and that the values on which meaning depends are neither simply invented nor simply discovered.

\paragraph{Nozick.} \citet[ch.~6]{Nozick1981} connected meaning with transcending limits: a life acquires meaning by connecting with value beyond itself, and the question of the meaning of that larger value recurs until it ends in something unlimited. He also gave, in the experience machine \citep[42--45]{Nozick1974}, the case that ETH-05 uses to show that well-being includes contact with reality and that EXI-17-a uses to show that meaning requires engagement with a real object (its Principle~P2).

\paragraph{Wolf.} \citet{Wolf1997} argued that meaning is a value distinct from happiness and from morality, and that a life is meaningful insofar as its subject is actively engaged in projects of worth. Her Tanner Lectures \citep{Wolf2010} developed the view, defended it against subjectivism and against purely objective accounts, and argued that the objective element should be understood as worth that does not derive from the subject's attitudes, while declining to give a general theory of that worth. They were published with comments by John Koethe, Robert M. Adams, Nomy Arpaly and Jonathan Haidt, and Wolf's replies. Wolf's view is the source of Definition~1.

\paragraph{Metz.} \citet{Metz2013} surveyed and classified the theories, argued against supernaturalism and subjectivism, and defended an objective naturalist theory, the fundamentality theory, on which a life is more meaningful the more its subject positively orients her rational nature toward conditions that are responsible for much else about human life. Metz is cited in EXI-01-a for the claim that there are projects of objective worth (\Lm{3.5}).

\paragraph{Later work.} The last fifteen years have added work on the temporal structure of meaning \citep{Kauppinen2012}, on the relation between meaning and consequences \citep{Smuts2013}, on cosmic insignificance \citep{Kahane2014}, on the pessimist's case that meaning is limited at best \citep{Benatar2017}, on the dependence of the meaning of our lives on the survival of humanity after us \citep{Scheffler2013}, on meaning in imperfect lives \citep{Landau2017} and on the variety of meaningful lives \citep{May2015}. Each enters the Commentary at the lemma to which it is relevant.

\subsection{The form of EXI-01-a}

EXI-01-a is written as a small axiomatic theory. It contains a definition (Definition~1), two assumptions (Assumptions~2 and~3), a theorem (Theorem~4) with a proof, and a remark (Remark~5) that says what the theorem answers and what it does not. The numbering is continuous from 1 to 5 and is not divided by section, as the numbering of longer manuscripts of the discipline is. Three features of the form matter for interpretation.

First, the form separates what is stipulated from what is asserted. Definition~1 is not true or false in the way an assumption is; it says what the text means by \enquote{meaningful to degree $d$}. Whether it captures what we ordinarily mean by a meaningful life is a question about the definition, and the definition is cited to Wolf because Wolf argued that it does. Assumptions~2 and~3 are assertions, and the text does not argue for them: the first is cited to Metz, and the second is left to its consequences and to the other manuscripts of its family (EXI-01-c, EXI-17-a). The theorem follows from the definition and the assumptions by a proof that occupies three sentences. Everything philosophically contested is therefore located in one of four places: the adequacy of the definition, the truth of either assumption, or the bridge between the assumptions and the theorem. The Commentary takes these in order.

Second, the form makes the text unusually easy to formalize, and it is formalized: Theorem~4 is proved in a Lean module whose axioms are listed in Section~\ref{s:lean}. The formalization does what formalization does in this discipline: it shows that the conclusion follows, and it shows exactly which premises it follows from. The second service is the more interesting here. It shows that Assumption~2 is not among the premises, and that a premise not stated in the text is (Propositions~\ref{p:redundant} and~\ref{p:bridge}).

Third, the brevity of the form leaves much unsaid that a longer text would say: the defence of the assumptions, the replies to objections, the relation to rival views. Remark~5 acknowledges this and refers the reader to the Commentary for two of the missing parts. The other missing parts are supplied in part by the other manuscripts of the family. EXI-01-c gives a second proof of Theorem~4 that does not use Assumption~3. EXI-01-d draws out what the availability claim implies and what it does not. The Commentary draws on both.

\subsection{The text in its discipline}

The discipline to which EXI-01-a belongs has twenty families, and most of them bear on it. We list the relations here so that later comments can refer to them briefly; Appendix~\ref{app:xref} gives the full table, computed from the numbered statements of every manuscript.

EXI-17 holds that a meaningful life requires active engagement with something worthwhile and that achievement is not necessary. Its principal, EXI-17-a, defines active engagement (Definition~2.1) and worth (Definition~2.2) in terms that agree with EXI-01-a, states the fitting-fulfilment view as its Principle~P1, adds a principle of contact with reality (Principle~P2), and derives that meaning is not hostage to luck in outcomes (Corollary~3.6), which supports the availability claim of EXI-01. It is the manuscript closest to EXI-01-a, and the two should be read together.

EXI-04 holds that life is absurd in Nagel's sense but that the appropriate response is irony rather than despair. EXI-04-a accepts the conclusion of EXI-01 and argues that the absurd survives it (its Objection~2 and reply), and it cites Remark~5 of EXI-01-a. EXI-04-b reconsiders Camus's Sisyphus. EXI-06 holds that the universe is indifferent and that indifference is not hostility, and EXI-06-b argues that this indifference takes nothing from meaning as EXI-01 understands it. EXI-11 holds that meaning does not require God, though a theistic framework would add to it; its principal defines meaning as EXI-01-a does (its Definition~2.1) and states as a principle that worth is not conferred by assignment (its Principle~P1). EXI-03 holds that an endless life would exhaust the categorical desires that give life its point, and EXI-03-a identifies those desires, in line with EXI-01, with desires for projects the person takes to be worthwhile. EXI-19 holds that death has no meaning in itself but gives life a shape by limiting it, and EXI-19-b draws out what this implies for meaning as EXI-01 and EXI-17 understand it. EXI-07 holds that happiness is a by-product of other pursuits; EXI-09 that existence precedes essence for human beings; EXI-20 that whether producing a body of work was meaningful depends on whether it was engagement with something worthwhile. Each of these is relevant at one or more lemmata.

Outside the discipline, the text depends on ETH-01, which holds that the good is what idealized reflective endorsement converges on and that it has seven irreducible components, and on the metaethical families MTE-01, MTE-05, MTE-09 and MTE-12, which hold that there are moral facts and that the arguments from queerness, evolution and disagreement do not show otherwise. These bear on Assumption~2. It touches ETH-05, which holds that well-being includes contact with reality, and REL-01, which holds that a maximally great being exists; the second bears on the additional meaning that EXI-11 assigns to a theistic framework, and so on the degree but not the availability of meaning.

\subsection{Note on the text}

The text of EXI-01-a consists of a title and eleven sentences, which contain \cnTokens{} words when citations are not counted, followed by a bibliography of three items. Its statements are introduced by bold labels (\enquote{Definition~1.}, \enquote{Assumption~2.} and so on) rather than by the numbered environments of the longer manuscripts, and the proof of Theorem~4 is introduced by an italic label and closed by a square. The Lean module refers to the statements by the same numbers: its docstrings name \enquote{Definition~1, positive case}, \enquote{Assumption~3} and \enquote{Theorem~4}, and they agree with the text. The text uses two British spellings, \enquote{defence} and \enquote{sceptical}, against the American spelling of the rest of the discipline; we quote them as they stand. The last sentence of Remark~5 refers to the present work by its identifier. No variant of the text is known to us.

\subsection{The commentary form and its precedent}

The commentary on an axiomatic text has a long precedent, and the precedent illuminates EXI-01-a. Proclus's commentary on the first book of Euclid's \emph{Elements} \citep{Proclus1970} comments on the definitions, the postulates, the common notions and the propositions in order, and distinguishes their roles: a definition says what is meant, a postulate demands something to be granted that is neither self-evident nor proved, a common notion states what is evident to all, and a proposition is demonstrated. Much of what he says about the postulates concerns their status. He held that the fifth postulate, on parallels, did not belong among the postulates, since it was not evident and its converse was proved as a proposition, and that it should be proved as a theorem; and he offered a proof, which later readers found to rely on an assumption equivalent to the postulate itself.

EXI-01-a has the same parts in miniature. Definition~1 says what is meant. Assumption~2 is closer to a common notion in the intention of its author, who cites rather than argues for it, and the Commentary has argued that its role is interpretive rather than demonstrative (Proposition~\ref{p:redundant}). Assumption~3 is a postulate in Proclus's sense: it carries the weight of the theorem, it is not self-evident, and it is the premise a critic will press. The family's response to it is Proclus's response to the fifth postulate. EXI-01-c proves it as a theorem from a premise thought to be better supported, its Assumption~2.2 (Reach), and anticipates the objection that Proclus's proof met, that the new premise is the old one in other words; its reply is that the new premise is supported by the ubiquity of persons and of the world as an object of understanding, and that the cases that threatened Assumption~3 satisfy it. Whether the reply succeeds is discussed at \Lm{4}. The parallel shows at least that the question it raises is the right one: in an axiomatic text, the dialectic concentrates on the premise that is neither definitional nor evident, and the commentary's task is to say whether it can be derived, from what, and at what cost.

\subsection{Method of this volume}

The text is divided in Section~\ref{s:text} into \cnSentLemmata{} sentence lemmata, numbered \Lm{1} (the title) to \Lm{12}, and \cnPhraseLemmata{} phrase lemmata, numbered by sentence and phrase, for the sentences that carry the argument: the definition, the two assumptions, the theorem, the central step of the proof and the remark on the absurd. Every word of each of those sentences belongs to at least one phrase lemma, except connectives whose role is discussed under the sentence lemma.

Sections~\ref{s:title} to~\ref{s:remark} comment on the lemmata in order. Each section opens with the sentence lemma, comments on the sentence as a whole, and then takes its phrases in turn. Where a phrase raises a question that the Commentary can settle, we state the result as a numbered proposition and prove it; where it raises a question that belongs to a later volume, we say so. Section~\ref{s:bib} comments on the bibliography of EXI-01-a, Section~\ref{s:lean} on the Lean module, declaration by declaration, and Section~\ref{s:whole} on the argument as a whole: the readings it admits, the objections to it that are not objections to any one lemma, and what it does not establish.

The appendices are computed. Appendix~\ref{app:conc} is a concordance: every word type of EXI-01-a, with the lemmata in which it occurs and its occurrences in each of the other \cnOtherMss{} manuscripts of the discipline, a table of fifteen term families by manuscript, and a keyword-in-context listing of every occurrence of the content words of the text in those manuscripts. Appendix~\ref{app:xref} lists every numbered statement of those manuscripts, read from their compiled \texttt{.aux} files, with its concept classes and the lemmata of EXI-01-a it bears on by a stated rule. Appendix~\ref{app:test} gives the testimonia. Each appendix has a preface stating what was computed and how.
